\documentclass[10pt]{article}
\usepackage[margin=0.8in]{geometry}
\usepackage{amsmath,booktabs,hyperref,xurl}
\usepackage[T1]{fontenc}
\setlength{\parindent}{0pt}
\setlength{\parskip}{5pt}
\title{Increase the configured collision margin to 0.20 m}
\author{Collision Avoidance Research}
\date{October 3, 2026}
\begin{document}
\maketitle
\textbf{Change.} The default \texttt{collision.safetyMarginMeters} is increased
from 0.10 to 0.20 m. This changes the existing sampled collision rows and
terminal separation reserve. All other controller parameters, the 10-m path
deviation bound, one quadratic CLF, shifted input initialization, and potential-field
restart logic are unchanged. No interval-clearance correction is introduced.
The study analysis description now refers to the configured buffer instead of
hard-coding the former value.

\textbf{Validation.} All 411 controller-related tests across 19 suites pass.
The initial run passed 410 tests; a terminal test used a fixed lower bound tied
to the former margin. Its expected geometric reserve now subtracts the configured
margin, and the complete 28-test terminal suite was rerun successfully.
The other 383 passing results remain from the unchanged all-suite run.
The scenario campaign started after tests completed, with MATLAB R2026a,
a single computational thread, exact ego observations and known constant target
acceleration/sideslip. Seven prescribed collision threats run at 8 and 15 m/s,
50-ms holds, 8/16 prefix steps and an 80-step completion. ODE45 replay uses
$10^{-11}/10^{-12}$ tolerances and 31 samples per hold. No random seed or
calculation delay is injected. Success requires positive sampled clearance,
no interruption, the 10-m limit, and a one-second post-conflict nominal dwell.
The 1500-hold observation cap is not a recovery deadline.

\textbf{Outcome.} 13/14 scenarios meet these criteria, with 1 solver interruptions, 0 sampled collision cases and 0 path-limit violations. Across 3896 holds, minimum sampled clearance is 14.597 cm (previous campaign: 5.849 cm); maximum path offset is 7.797 m.

\begin{center}
\begin{tabular}{lrrrr}
\toprule
Speed / scenario & Old gap (cm) & New gap (cm) & Peak offset (m) & Recovery (s) \\
\midrule
8 / headOn & 139.28 & 139.28 & 5.453 & 10.75 \\
8 / acceleratingHeadOn & 106.67 & 106.67 & 5.314 & 10.70 \\
8 / brakingLead & 24.57 & 22.17 & 5.381 & 15.15 \\
8 / crossing & 9.24 & 14.60 & 5.667 & 10.80 \\
8 / turningCrossing & 5.85 & 19.27 & 5.597 & 10.40 \\
8 / curvedHeadOn & 178.90 & 178.90 & 5.609 & 9.70 \\
8 / curvedCrossing & 9.35 & 18.64 & 5.776 & 9.75 \\
15 / headOn & 21.56 & 21.56 & 5.093 & 8.95 \\
15 / acceleratingHeadOn & 8.47 & 1344.32 & 0.133 & -- \\
15 / brakingLead & 25.16 & 26.24 & 7.797 & 13.55 \\
15 / crossing & 135.80 & 135.82 & 6.557 & 10.60 \\
15 / turningCrossing & 129.99 & 129.98 & 6.514 & 9.50 \\
15 / curvedHeadOn & 191.87 & 191.44 & 6.209 & 64.75 \\
15 / curvedCrossing & 100.69 & 100.69 & 6.365 & 9.95 \\
\bottomrule
\end{tabular}
\end{center}
\textbf{Measured runtime.} Whole successfully returned running-frame median/maximum: 23.862/93.064 ms. Largest first call: 238.656 ms. There are 0 running calls above 100 ms. These are shared-workstation measurements without an injected execution delay.

\textbf{Scope.} 3 scenes still have a sampled physical gap below the configured 20 cm. The larger node buffer is an empirical clearance reserve; it does not close the intersample and approximation gap diagnosed in \path{report/MARGIN_SAMPLING_DIAGNOSIS_20261003.tex}. Neither a continuous-time 20-cm guarantee nor a real-time deadline guarantee is asserted.

\textbf{Interrupted case.} At 15 m/s, accelerating head-on stops before
issuing the sixth control, at 0.25 s. The 13.443-m minimum in its table row
covers only five executed holds and does not demonstrate avoidance of that
encounter. Direct replay of the frozen failed frame reproduces
\texttt{linearizationAccuracyNotReached}: the shifted model's collision-relevant
pose disagreement is 10.535 mm against a 10-mm threshold. Its anchor has an
inequality residual of 0.041818, so feasibility-preserving damping has no
accepted base point. Fresh potential-field initialization permits both convex
solves but the full step has 26.947-mm disagreement; three damped trials still
have 19.701, 18.003 and 17.579 mm. Thus finite convex solutions exist, but neither
allowed anchor supplies a control meeting the unchanged model-accuracy gate.
No control is issued by this failed call. The separate replay runs in a fresh
MATLAB process and its elapsed time is not a warm runtime measurement. The
campaign's returned-frame statistics above exclude the unrecorded failed-call
latency. This is a regression relative to the 0.10-m campaign, not evidence of
physical impossibility. No relaxation of acceptance or other controller tuning
is included in this parameter change.

\textbf{Artifacts.} Run \texttt{runPotentialFieldCampaign}; compact metrics,
source hashes, original-result provenance and test records are under
\path{report/MARGIN_20CM_20261003/}. The baseline is the 0.10-m path-bound campaign
at commit \path{8347c2aac7952ae8065f9c571e3c46cdca77e049}.
\end{document}
